\uuid{6cAX}
\exo7id{4433}
\titre{exo7 4433}
\auteur{quercia}
\organisation{exo7}
\datecreate{2010-03-14}
\isIndication{false}
\isCorrection{true}
\chapitre{Série numérique}
\sousChapitre{Autre}
\cpge{mpsi-mp2i.procedes-sommatoires-discrets.a, pcsi.series-numeriques.a, ptsi.series-numeriques.a, mpsi-mp2i.calcul-algebrique-trigonometrie.a, pcsi.calcul-algebrique-trigonometrie.a, ptsi.calcul-algebrique-trigonometrie.a, bcpst1.methodes-calcul, mp-mpi.series-numeriques-vectorielles.a, pt.series-numeriques.b, bcpst2.series-reelles}
\module{Analyse}
\niveau{L1}
\difficulte{}

\contenu{
\question{
On admet que $\sum_{k=1}^\infty \frac1{k^2} = \frac{\pi^2}6$.
Calculer $\sum_{k=1}^\infty \frac 1{k^2(k+1)^2}$.
}
\reponse{
$\frac{\pi^2}3 - 3$.
}
}
